\textbf{Proof.}
For arm \(a\), write \(F_{a,m_a}\) for the distribution function of \(\widehat F_a\).  We first prove the required empirical-law convergence directly.  For fixed rational \(t\in[0,1]\), the variables
\[
 X_{j,a}(t)=\mathbf 1\{Y^P_{j,a}\le t\}
\]
are independent Bernoulli variables with mean \(F_a(t)\) by \(\texttt{ass:unpaired-pilot}\).  If \(X\) is Bernoulli with mean \(p\), then
\[
 \mathbb E\exp\{\lambda(X-p)\}
 =e^{-\lambda p}(1-p+pe^\lambda)\le e^{\lambda^2/8}.
\tag{1}
\]
Indeed, the second derivative in \(\lambda\) of \(\log(1-p+pe^\lambda)\) is the variance of a tilted Bernoulli variable and is at most \(1/4\), while that logarithm minus \(p\lambda\) and its first derivative vanish at zero.  Exponential Markov inequality, independence, and optimization at \(\lambda=4s\) therefore give, for \(s>0\),
\[
 P\!\left(\left|F_{a,m_a}(t)-F_a(t)\right|>s\right)
 \le 2e^{-2m_as^2}.
\tag{2}
\]
For each rational \(s>0\), the right-hand side is summable in \(m_a\).  The union bound yields
\[
 P\!\left(\limsup_{m_a\to\infty}
 \{|F_{a,m_a}(t)-F_a(t)|>s\}\right)
 \le \lim_{N\to\infty}\sum_{m_a\ge N}2e^{-2m_as^2}=0.
\tag{3}
\]
Intersecting the probability-one events in (3) over the countable rational pairs \((t,s)\) proves simultaneous convergence at every rational \(t\).

Fix an integer \(K\ge1\) and the rational grid \(x_j=j/K\).  Monotonicity of the two distribution functions implies, for \(y\in[x_j,x_{j+1}]\),
\[
 |F_{a,m_a}(y)-F_a(y)|
 \le
 \max_{r\in\{j,j+1\}}|F_{a,m_a}(x_r)-F_a(x_r)|
 +F_a(x_{j+1})-F_a(x_j).
\tag{4}
\]
After integration and summation over \(j\), pointwise convergence on this finite grid and telescoping of the increments in (4) give
\[
 \limsup_{m_a\to\infty}
 \int_0^1|F_{a,m_a}(y)-F_a(y)|\,dy\le K^{-1}
 \quad\text{almost surely}.
\tag{5}
\]
Letting \(K\to\infty\) proves that the integral in (5) converges to zero.

For any distribution functions \(R,S\) on \([0,1]\), Tonelli's identity
\[
 \int_0^1|R^{-1}(u)-S^{-1}(u)|\,du
 =\int_0^1|R(y)-S(y)|\,dy
\tag{6}
\]
follows by writing the absolute difference of two numbers as the integral of the absolute difference of their threshold indicators and using
\(
 \mathbf 1\{R^{-1}(u)\le y\}=\mathbf 1\{u\le R(y)\}
\)
outside a null set.  Equations (5)--(6) show
\[
 \int_0^1|\widehat F_a^{-1}(u)-F_a^{-1}(u)|\,du\longrightarrow0
 \quad\text{almost surely},\qquad a\in\{1,0\}.
\tag{7}
\]
The bounded support required here is exactly \(\texttt{ass:bounded-outcomes}\).

Because \(y\mapsto y\) and \(y\mapsto y^2\) are respectively one- and two-Lipschitz on \([0,1]\), (7) gives convergence of the empirical means, second moments, and variances.  Moreover, for \(x,x',y,y'\in[0,1]\),
\[
 |xy-x'y'|\le |x-x'|+|y-y'|.
\tag{8}
\]
Applying (8) to the comonotone quantiles and to the antimonotone quantiles proves almost-sure convergence of both empirical cross moments and hence of both empirical covariance endpoints.

By \(\texttt{def:moment-matrix}\) and the endpoint-matrix construction in \(\texttt{def:pilot-plugin}\), every block of \(\widehat Q_e-Q(c_e)\) is a convergent scalar multiplying either \(I_n\) or \(\mathbf 1\mathbf 1^{\top}\).  Since \(n\) is fixed,
\(
 \lVert I_n\rVert_{\mathrm{op}}=1
\)
and
\(
 \lVert\mathbf 1\mathbf 1^{\top}\rVert_{\mathrm{op}}=n
\).
Therefore
\[
 M_m=\max_{e\in\{L,U\}}
 \lVert\widehat Q_e-Q(c_e)\rVert_{\mathrm{op}}
 \longrightarrow0
 \quad\text{almost surely under }P.
\]
The countable first-qualifying basis selector in \(\texttt{def:pilot-plugin}\) is not used in this endpoint calculation. \(\square\)